% sec_04_fidelity_grid
\section{The fidelity grid}\label{sec:grid}

AEGIS exposes nine fidelity levels, numbered 0 to 8, that trade accuracy for
cost. Rather than nine unrelated models, they are organised as an orthogonal
composition: a small set of physical corrections, each touching a different
structural element of the master law, switched on independently. This section
states the composition, lists the levels, and identifies the new orthogonal
axes (diffraction model and inter-body coupling) developed later in the report.

\subsection{Two regimes}\label{sec:two-regimes}

The levels split into two regimes by what they act on.

The incoherent levels (0 to 6) act on per-path power. Each propagation path
contributes a scalar power density $S_i$, and absorbed power densities add as
powers,
\begin{equation}\label{eq:incoherent-sum}
  \Sab(\rr) \;=\; \sum_i S_i\;\Teff^{(i)}(\rr)\;g\!\bigl(\mu_i(\rr),\sigma\bigr)\;
  O_i(\rr,\khat_i),
\end{equation}
the multi-path form of the master law (\cref{eq:master-law-exact}) with the
activation $g$ defined below. Phase information between paths is discarded.

The coherent levels (7 and 8) act on complex amplitudes. Each path carries a
polarisation-amplitude vector $\psivec_i$, and the body-surface field is the
coherent sum over paths driven by a precoder. The absorbed power density becomes
a Hermitian quadratic form in the precoder $\xx$,
\begin{equation}\label{eq:coherent-form}
  \Sab(\rr) \;=\; \norm{\Gtil(\rr)\,\xx}^2
  \;=\; \xx\herm\,\Qop(\rr)\,\xx,
\end{equation}
with $\Gtil$ the body-surface channel and $\Qop$ the exposure operator, both
developed in \cref{sec:coherent}. \Cref{eq:coherent-form} reduces to
\cref{eq:incoherent-sum} when path phases are randomised.

\subsection{Orthogonal corrections}\label{sec:orthogonal}

Within the incoherent regime the refinements compose because each modifies a
different factor of the per-path term, as shown in
\texttt{composability\_analysis.md}. The general combined per-triangle law is
\begin{equation}\label{eq:combined}
  \Sab \;=\; \underbrace{\Bigl[\Tavg(\mu)+\tfrac{q}{2}\Delta T(\mu)\Bigr]}_{\text{Fresnel}\,+\,\text{polarisation}}
  \;\underbrace{g(\mu,\sigma)}_{\text{activation}}\;S
  \;+\;\underbrace{T_0\,\frac{H}{k}\,g(\mu,\sigma)^2\,S}_{\text{curvature}},
\end{equation}
where the axes are:
\begin{itemize}[leftmargin=1.4em,itemsep=1pt]
  \item \textbf{Fresnel} ($T_0\to\Tavg(\mu)$): the angle-dependent transmission
    of \cref{sec:fresnel}. Modifies the material factor.
  \item \textbf{Polarisation} (the $\tfrac{q}{2}\Delta T$ term): restores the
    TM excess of \cref{sec:pol}, with $q$ the local TM excess. Modifies the
    material factor only.
  \item \textbf{Curvature} (the $H/k$ term): the first-order physical-optics
    correction for a surface of twice-mean-curvature $H$. Adds a perturbative
    $g^2$ term.
  \item \textbf{Diffraction} ($g\in\{\ReLU,\,\mathrm{GeLU},\,\text{Fock}\}$):
    replaces the hard back-face gate $\ReLU(\mu)$ with a physically derived
    smooth transition near the shadow boundary. Modifies the activation only.
  \item \textbf{Inter-body} (the visibility factor $O$): reflections and
    occlusion from other bodies in the scene. Modifies the visibility factor.
\end{itemize}
The corrections are structurally orthogonal: no two modify the same factor, so
any subset can be enabled and the cross-terms between corrections are smaller
than the corrections themselves (for example the polarisation-curvature cross
term is $\sim0.03\%$ for an arm at \SI{28}{\giga\hertz}, three orders of
magnitude below the leading curvature term).

\subsection{The new orthogonal axes}\label{sec:new-axes}

Two axes generalise the original six-level ladder and are the subject of the
later sections. The \textbf{diffraction model} is promoted from a binary
on/off switch to a choice $g\in\{\text{none},\,\text{gelu},\,\text{fock}\}$:
\texttt{none} is the hard $\ReLU$, \texttt{gelu} is the curvature-scaled error
function of \cref{eq:combined}, and \texttt{fock} is the rigorous Fock
creeping-wave model (\cref{sec:fock-local,sec:cmetric}). The
\textbf{inter-body} axis turns the visibility factor $O$ into a coupling that
carries reflected and occluded power between bodies (\cref{sec:interbody}).
Both axes are orthogonal to the Fresnel, polarisation, and curvature axes and
compose with them.

\subsection{Level map}\label{sec:level-map}

\Cref{tab:levels} maps each level to the corrections it enables. Levels 0 and 1
collapse the spatial map to a single aggregate number via the projected area and
absorption directivity. (The viewer's Level 0/1 implementation carries a
documented factor-of-4 caveat in the area normalisation $\Aab$, traced in
\texttt{composability\_analysis.md}: the Cauchy relation
$\langle\Aperp\rangle=A_{\mathrm{total}}/4$ is applied once too often, leaving
the aggregate power four times too small unless $\Aab=A_{\mathrm{total}}$.)

\begin{table}[!htbp]
\centering
\caption{The nine fidelity levels as an orthogonal composition. Fresnel:
  angle-dependent $\Tavg$. Pol: TM excess $\tfrac{q}{2}\Delta T$. Curv:
  physical-optics $H/k$ term. Diff: activation $g$. Coh: coherent quadratic
  form.}
\label{tab:levels}
\begin{tabular}{clccccl}
\toprule
Level & Name & Fresnel & Pol & Curv & Diff & Output \\
\midrule
0 & Worst-case bound      & $T_0$    &   &   & ReLU & scalar bound \\
1 & Aggregate directivity & $T_0$    &   &   & ReLU & scalar power \\
2 & Geometric map         & $T_0$    &   &   & ReLU & $\Sab(\rr)$ \\
3 & Fresnel map           & $\Tavg$  &   &   & ReLU & $\Sab(\rr)$ \\
4 & Polarisation          & $\Tavg$  & \checkmark &   & ReLU & $\Sab(\rr)$ \\
5 & Curvature             & $\Tavg$  & \checkmark & \checkmark & ReLU & $\Sab(\rr)$ \\
6 & Diffraction           & $\Tavg$  & \checkmark & \checkmark & GeLU/Fock & $\Sab(\rr)$ \\
7 & Coherent              & exact $t_s,t_p$ & intrinsic & opt. & opt. & $\norm{\Gtil\xx}^2$ \\
8 & ECBF                  & exact $t_s,t_p$ & intrinsic & opt. & opt. & $\xx\herm\Qop\xx$ \\
\bottomrule
\end{tabular}
\end{table}

In the coherent levels (7 and 8) polarisation is intrinsic: the channel applies
the complex amplitude coefficients $t_s,t_p$ per TE/TM component, so no separate
polarisation switch is needed, and curvature and diffraction enter as real
per-(triangle, path) weights that preserve the quadratic-in-$\xx$ structure.
The remaining sections derive each axis in turn.
